\documentclass[12pt, a4paper]{article}
\usepackage[T1]{fontenc}
\usepackage{lmodern}
\usepackage{amsmath, amsthm, amssymb, mathtools}
\usepackage[margin=1in]{geometry}
\usepackage{xr-hyper}
\usepackage[colorlinks=true,linkcolor=blue,citecolor=blue]{hyperref}
\usepackage{cleveref}
\usepackage{graphicx, subcaption}
\usepackage{natbib}
\usepackage{booktabs}
\usepackage{enumerate}

% Code references: scaled monospace with line-breaking
\newcommand{\code}[1]{\texttt{\small #1}}

\newtheorem{theorem}{Theorem}[section]
\newtheorem{lemma}[theorem]{Lemma}
\newtheorem{proposition}[theorem]{Proposition}
\newtheorem{corollary}[theorem]{Corollary}
\newtheorem{definition}{Definition}[section]
\newtheorem{remark}{Remark}[section]
\newtheorem{conjecture}[theorem]{Conjecture}
\newtheorem{derivation}[theorem]{Derivation}
\newtheorem{observation}[theorem]{Observation}
\numberwithin{equation}{section}

\newcommand{\dd}{\mathrm{d}}
\newcommand{\seriestitle}{Why Rotating Fluids Make Polygons}


\title{\seriestitle\\[6pt]\large A Guide for Readers}
\author{Gordon Speagle}
\date{}

\begin{document}
\maketitle

%% ====================================================================
\section{What this series claims}
%% ====================================================================

The stability eigenvalue of a regular polygon of point vortices
splits into two pieces:
\begin{equation}\label{eq:central}
  \lambda_m = C_1(\text{surface}) - \frac{m(N{-}m)}{2}.
\end{equation}
The first piece depends on the surface geometry.
The second depends only on the number of vortices~$N$ and the
mode index~$m$.
This one formula, applied on the Seifert manifold
$\mathbf{H}^2 \times_N S^1$, produces
the Standard Model gauge group, the Weinberg angle,
the CKM phase, three generations of fermions,
the cosmological constant, the Hubble parameter,
and the dark matter fraction---from one integer ($N = 7$)
and one scale ($M_P$).

The series is six papers.
Papers~I--III develop the mathematics, the vortex physics,
and the gravitational theory.
Paper~IV derives gauge theory and particle physics from $N = 7$.
Paper~V derives cosmology from $N = 11$.
Paper~VI collects the predictions, parameter accounting,
and status classification.

%% ====================================================================
\section{What is new}
%% ====================================================================

\subsection{New mathematics}

Results proved for the first time in this series.

\medskip
\begin{tabular}{p{6.5cm}lp{4cm}}
\toprule
Result & Equation / Ref & Paper \\
\midrule
Sign rule: $H''(0) < 0$ for any
  $h$ with $h' < 0$, $(rh')' \le 0$ &
  Thm~I-2.1 & I \\[4pt]
Constrained energy minimum: the $N$-gon
  minimises $H$ on $\{L = \mathrm{const}\}$
  for $N \le 7$ &
  Thm~I-2.2 & I \\[4pt]
Metric-independence of $f(m,N) = m(N{-}m)/2$
  on all constant-curvature surfaces &
  Thm~I-3.2 & I \\[4pt]
$\mathbf{S}^2$ thresholds are all rational:
  $\xi_{\mathrm{crit}} \in \{1/3, 1/5, 1/7, 1/19\}$ &
  eq.~I-(4.4) & I \\[4pt]
$\mathbf{H}^2$ thresholds are Pell units in
  quadratic fields &
  Prop~I-5.1 & I \\[4pt]
Algebraic-integer classification:
  $N \in \{8, 9, 11, 15, 23\}$ only &
  Prop~I-5.1 & I \\[4pt]
Field extension pattern:
  even $N \to \mathbb{Q}(\sqrt{\mathrm{sq\_free}(N{-}1)})$,
  odd $N \to \mathbb{Q}(\sqrt{\mathrm{sq\_free}(N{-}3)})$ &
  eq.~I-(5.1) & I \\[4pt]
Geodesic correspondence: five thresholds
  $=$ five closed geodesics on the modular surface,
  traces $T = 2 + 16/(N{-}7)$ &
  Prop~I-5.9 & I \\[4pt]
Universality: $N_{\mathrm{crit}} = 7$ at $\xi = 0$
  on every quotient $\Gamma\backslash\mathbf{H}^2$ &
  Thm~I-5.12 & I \\[4pt]
Higher-genus palindromic hierarchy:
  degree-$2d$ palindromic polynomial for
  trace-field degree $d$ &
  Conj~I-5.14 & I \\[4pt]
Quadratic reciprocity as a selection rule
  for visible thresholds on $\Gamma_0(p)$ &
  Prop~I-5.10 & I \\[4pt]
Quartic normal form: $\alpha_0 = 45/14$
  resolves $N = 7$ marginal stability &
  \S II-5 & II \\[4pt]
CMS--CS Casimir identification:
  $C_2^{\mathrm{CMS}} = C_2^{\mathrm{CS}}$ via
  symmetric-space Laplacian &
  Prop~III-2.1 & III \\[4pt]
Polygon entropy maximisation implies
  $R = \mathrm{const}$ (Einstein equations) &
  Thm~III-6.1 & III \\[4pt]
\bottomrule
\end{tabular}

\subsection{New physics}

Physical results and predictions not previously in the literature.

\medskip
\begin{tabular}{p{6.5cm}lp{4cm}}
\toprule
Result & Ref & Paper \\
\midrule
The $N$-gon is a constrained energy
  \emph{minimum}, not a maximum
  (corrects the statistical-mechanics literature) &
  Thm~I-2.2 & I, II \\[4pt]
Blob bridge: finite-core vortex patches
  converge to point-vortex stability at
  $O(\varepsilon^2)$, profile-independent &
  Prop~II-7.1 & II \\[4pt]
Circulation disorder resilience:
  $N = 6$ stability domain is robust,
  not a knife-edge &
  \S II-9 & II \\[4pt]
Graviton emerges at $N = 7$ as a boundary
  Virasoro mode ($j = 2$, spin-$2$) &
  Prop~III-2.2 & III \\[4pt]
Vacuum Einstein equations from polygon
  entropy maximisation + Onsager selection &
  Thm~III-6.1 & III \\[4pt]
$\mathrm{SU}(3) \times \mathrm{SU}(2) \times
  \mathrm{U}(1)$ from $N = 7$ orbifold +
  McKay + CS gravity &
  Thms~IV-6.1, IV-6.3 & IV \\[4pt]
Weinberg angle $\sin^2\theta_W = 3/11$
  from CS conformal weight ratio &
  Thm~IV-9.1 & IV \\[4pt]
CKM phase $\delta = 70.2^\circ$ from
  BF-crossing scattering phase on
  $\mathbf{H}^2$ &
  Thm~IV-14.1 & IV \\[4pt]
$N = 11$ uniquely selected by KK flux
  additivity; $H_0 = 67.4$~km/s/Mpc &
  Thm~V-3.1 & V \\[4pt]
Dark matter from frozen KK polygon modes
  ($50$--$950$~TeV, absolutely stable) &
  \S V-5 & V \\[4pt]
Baryogenesis: $\eta_B \approx 2 \times 10^{-9}$,
  $v_w$-independent &
  \S V-7 & V \\[4pt]
\bottomrule
\end{tabular}

\subsection{New connections}

Known results connected for the first time in this series.

\medskip
\begin{tabular}{p{6.5cm}lp{4cm}}
\toprule
Connection & Ref & Paper \\
\midrule
Havelock eigenvalues $=$ CMS Casimir values
  $=$ CS Wilson-line Casimirs (all from the
  same $\mathrm{sl}(2,\mathbb{R})$ Laplacian) &
  Prop~III-2.1 & III \\[4pt]
Pell equation governs both vortex stability
  thresholds and the graviton spin ($j = 2$) &
  Prop~III-2.2 & I, III \\[4pt]
Vortex stability thresholds $=$ closed
  geodesics on the modular surface &
  Prop~I-5.9 & I \\[4pt]
$\xi^* = 8 - 3\sqrt{7}$ is the inverse
  fundamental unit of $\mathbb{Z}[\sqrt{7}]$ &
  eq.~I-(4.9) & I \\[4pt]
Onsager negative-temperature condensation
  $=$ WDW quantum cosmology
  (Feynman--Kac identity) &
  \S III-13 & III \\[4pt]
McKay correspondence + DHVW orbifold
  $\to$ $\mathrm{SU}(3)_1$ gauge sector &
  Thm~IV-6.3 & IV \\[4pt]
\bottomrule
\end{tabular}


%% ====================================================================
\section{What is known}
%% ====================================================================

Established results used as building blocks (not claimed as new).

\medskip
\begin{tabular}{p{6.5cm}p{6.5cm}}
\toprule
Result & Source \\
\midrule
Havelock eigenvalue formula $\lambda_m = (N{-}1) - m(N{-}m)/2$ &
  Havelock (1931) \\[3pt]
Calogero--Moser--Sutherland integrability &
  Calogero (1971), Sutherland (1971) \\[3pt]
Chern--Simons / WZW correspondence &
  Witten (1989) \\[3pt]
McKay correspondence ($\mathbb{Z}/n \subset \mathrm{SU}(2)
  \to A_{n-1}$ Dynkin diagram) &
  McKay (1980) \\[3pt]
Brown--Henneaux central charge $c = 3\ell/(2G)$ &
  Brown--Henneaux (1986) \\[3pt]
Thurston geometrisation of 3-manifolds &
  Thurston (1982), Perelman (2003) \\[3pt]
Born--Oppenheimer approximation &
  Born--Oppenheimer (1927) \\[3pt]
Selberg trace formula &
  Selberg (1956) \\[3pt]
Kaluza--Klein dimensional reduction &
  Kaluza (1921), Klein (1926) \\[3pt]
\bottomrule
\end{tabular}

%% ====================================================================
\section{Where to find it}
%% ====================================================================

\paragraph{Planetary scientist.}
Start at Paper~II, \S6 (the Rossby bridge) and \S8
(deformation radius).
The key result: the polygon is stable for $N \le 7$
(Paper~I, Corollary~5.8).
Saturn's hexagon ($N = 6$) and Jupiter's octagon ($N = 8$
with a central vortex) are discussed in Paper~II, \S4 and \S8.
Skip Papers~III--VI.

\paragraph{Fluid dynamicist.}
Paper~I proves the sign rule and the Havelock identity.
Paper~II proves the constrained minimum (Theorem~3.2),
resolves $N = 7$ (the quartic normal form, \S5),
and develops the blob and Rossby bridges (\S7).
The curved-surface thresholds (Paper~I, \S4--5) are new.
Skip Papers~III--VI.

\paragraph{Mathematical physicist / number theorist.}
The algebraic number theory is in Paper~I, \S5:
palindromic Pell units, the field extension pattern,
the geodesic correspondence, quadratic reciprocity
as a selection rule, and the higher-genus conjecture.
The K-theory classification of stability phases is in
Paper~I, \S7.
The CMS--CS Casimir identification is in Paper~III, \S2.

\paragraph{Particle physicist.}
Paper~IV derives the gauge group (\S8), the Weinberg angle
(\S9), fermion masses and CKM mixing (\S14).
The complete $4$D Lagrangian is in \S12.
For the logical foundations, read Paper~III, \S12
(Einstein equations from polygon stability) and \S13
(Onsager--WDW correspondence).

\paragraph{Cosmologist.}
Paper~V derives $H_0$, $\Omega_\Lambda$, $\Omega_{\mathrm{DM}}$,
$\Omega_b$, the neutrino mass spectrum, and baryogenesis from
the $N = 11$ instanton.
For the logical foundation, read Paper~IV, \S2 (the Seifert
manifold) and Paper~III, \S7 (the WDW equation).
The full predictions table is in Paper~VI.

\paragraph{Skeptic / referee.}
Read Paper~VI first: the parameter accounting (\S2) lists
every input and structural choice; the status classification
(\S3) labels each result as Theorem, Derivation, or Conjecture.
The falsifiable predictions (\S7) are binary: normal hierarchy,
$w = -1$, no $0\nu\beta\beta$, proton stable,
$\theta_{\mathrm{QCD}} = 0$.
Five experiments will test these within a decade.

%% ====================================================================
\section{The derivation chain}
%% ====================================================================

\begin{center}
\small
\begin{tabular}{lll}
\toprule
Paper & Input & Output \\
\midrule
I (Mathematics) & $\mathbb{Z}_N$ representation theory &
  Sign rule, Havelock identity, $N_{\mathrm{crit}} = 7$, \\
  & & Pell units, geodesic correspondence \\[3pt]
II (Physics) & Paper I + physical constraints &
  Constrained minimum, quartic at $N = 7$, \\
  & & blob/Rossby bridges, BEC/planetary predictions \\[3pt]
III (Gravity) & Papers I--II + $2{+}1$D gravity &
  CMS--CS identification, graviton at $N = 7$, \\
  & & Einstein equations, Onsager--WDW \\[3pt]
IV (Gauge theory) & Paper III + Seifert $+$ $N = 7$ &
  $\mathrm{SU}(3){\times}\mathrm{SU}(2){\times}\mathrm{U}(1)$,
  $\sin^2\theta_W = 3/11$, \\
  & & CKM phase, fermion masses \\[3pt]
V (Cosmology) & Paper IV + $N = 11$ &
  $H_0 = 67.4$, $\Omega_\Lambda$, DM, baryogenesis, \\
  & & neutrinos, inflation \\[3pt]
VI (Discussion) & Papers I--V &
  Parameter accounting, status, predictions \\
\bottomrule
\end{tabular}
\end{center}

\noindent
Each paper depends only on the papers above it.
Paper~I is self-contained.

\bibliographystyle{plainnat}
\bibliography{../shared/refs}

\end{document}
